% Test fixture for scripts/check_citations.py: a line-numbered, double-spaced
% AMS-style manuscript with known citation errors planted in it. The expected
% flags are asserted in test_check_citations.py; keep the two in step.
\documentclass[12pt]{article}
\usepackage[T1]{fontenc}
\usepackage[utf8]{inputenc}
\usepackage[margin=1in]{geometry}
\usepackage{setspace}
\usepackage{lineno}
\doublespacing
\linenumbers
\setlength{\parindent}{2em}
\begin{document}

\section*{Abstract}
Mesoscale convective systems in the tropics are studied with satellite data.

\section{Introduction}
Tropical heat is carried to the upper troposphere by deep convective towers
(Malkus and Riehl, 1958; Houze 2004), organised into mesoscale systems.
Their convective and stratiform regions were described by Zipser (1977) and
Houze (1989, 2018). Stratiform rain fractions come from Schumacher and
Houze (2003, 2006), and the latent heating from Liu et al. (2021).
Global tracking is described by Feng et al. (2021a, b), and the tracked data set
of Feng et al. (2021) is used here. Land and ocean differ (Nesbitt et al. 2000;
Mohr and Zipser 1996). Anvil clouds were identified by Yuan et al. (2010) and
Elsaesser et al. (2022). The IPCC (IPCC 2021) assesses future changes, and
van der Linden et al. (2020) and Volonté et al. (2022) examine case studies,
as do García-Pérez and Smith (2019) and Hagos et al. (2013).
10 km grid spacing is used in all simulations described in Section 2 of this
paper, and the period is January 2020 to March 2021.

Cold pools organise tropical convection (Cortado 2013, hereafter C13). The
Radiative-Convective Equilibrium Model Intercomparison Project (hereafter RCEMIP)
is not a citation. As C13 showed, and as later work confirmed, C13 still holds.

\section{Acknowledgments}
Supported by contract DE-AC05-76RL01830.

\section*{References}
\begin{flushleft}
\setlength{\parindent}{-2em}\setlength{\leftskip}{2em}

Cortado, M., 2013: Cold pools and the organisation of tropical convection. \textit{J. Atmos. Sci.}, \textbf{70}, 100--120.

Elsaesser, G. S., R. Roca, T. Fiolleau, A. D. Del Genio, and J. Wu, 2022: A simple model for tropical convective cloud shield area growth and decay rates informed by geostationary IR, GPM, and Aqua/AIRS satellite data. \textit{J. Geophys. Res. Atmos.}, \textbf{127}, e2021JD035599.

Elsaesser, G. S., R. Roca, T. Fiolleau, A. D. Del Genio, and J. Wu, 2022: A simple model for tropical convective cloud shield area growth and decay rates informed by geostationary IR, GPM, and Aqua/AIRS satellite data. \textit{J. Geophys. Res. Atmos.}, \textbf{127}, e2021JD035599.

Feng, Z., and Coauthors, 2021a: A global high-resolution mesoscale convective system database using satellite-derived cloud tops, surface precipitation, and tracking. \textit{J. Geophys. Res. Atmos.}, \textbf{126}, e2020JD034202.

Feng, Z., F. Song, K. Sakaguchi, and L. R. Leung, 2021b: Evaluation of mesoscale convective systems in climate simulations: Methodological development and results from MPAS-CAM over the United States. \textit{J. Climate}, \textbf{34}, 2611--2633.

Hagoss, S. M., C. Zhang, W.-K. Tao, S. Lang, Y. N. Takayabu, S. Shige, M. Katsumata, B. Olson, and T. L'Ecuyer, 2013: Estimates of tropical diabatic heating profiles. \textit{J. Climate}, \textbf{23}, 542--558.

Houze, R. A., Jr., 2004: Mesoscale convective systems. \textit{Rev. Geophys.}, \textbf{42}, RG4003.

------, 2018: 100 years of research on mesoscale convective systems. \textit{A Century of Progress in Atmospheric and Related Sciences: Celebrating the American Meteorological Society Centennial, Meteor. Monogr.}, No. 59, Amer. Meteor. Soc., 17.1--17.54.

IPCC, 2021: \textit{Climate Change 2021: The Physical Science Basis}. Cambridge University Press, 2391 pp.

Liu, N., L. R. Leung, and Z. Feng, 2021: Global mesoscale convective system latent heating characteristics from GPM retrievals and an MCS tracking dataset. \textit{J. Climate}, \textbf{34}, 8599--8613.

Mapes, B. E., 1993: Gregarious tropical convection. \textit{J. Atmos. Sci.}, \textbf{50}, 2026--2037.

Mohr, K. I., and E. J. Zipser, 1996: Mesoscale convective systems defined by their 85-GHz ice scattering signature: Size and intensity comparison over tropical oceans and continents. \textit{Mon. Wea. Rev.}, \textbf{124}, 2417--2437.

Nesbitt, S. W., E. J. Zipser, and D. J. Cecil, 2000: A census of precipitation features in the tropics using TRMM: Radar, ice scattering, and lightning observations. \textit{J. Climate}, \textbf{13}, 4087--4106, https://doi.org/10.1175/1520-0442(2000)013,4087:ACOPFI.2.0.CO;2.

García Pérez, L., and J. Smith, 2019: A made-up entry whose surname has lost its hyphen. \textit{J. Climate}, \textbf{32}, 1--20.

Riehl, H., and J. S. Malkus, 1958: On the heat balance in the equatorial trough zone. \textit{Geophysica}, \textbf{6}, 503--538.

Schumacher, C., and R. A. Houze Jr., 2003: Stratiform rain in the tropics as seen by the TRMM Precipitation Radar. \textit{J. Climate}, \textbf{16}, 1739--1756.

------, and ------, 2006: Stratiform precipitation production over sub-Saharan Africa and the tropical East Atlantic as observed by TRMM. \textit{Quart. J. Roy. Meteor. Soc.}, \textbf{132}, 2235--2255.

van der Linden, R., A. H. Fink, J. G. Pinto, and T. Phan-Van, 2020: The dynamics of an extreme precipitation event in northeastern Vietnam in 2015. \textit{Mon. Wea. Rev.}, \textbf{148}, 2995--3015.

Volonté, A., A. J. Turner, R. Schiemann, P. L. Vidale, and N. P. Klingaman, 2022: The interaction of tropical and extratropical air masses controlled the 2018 Kerala floods. \textit{Quart. J. Roy. Meteor. Soc.}, \textbf{148}, 3011--3033.

Yuan, J., and R. A. Houze Jr., 2010: Global variability of mesoscale convective system anvil structure from A-Train satellite data. \textit{J. Climate}, \textbf{23}, 5864--5888.

Zipser, E. J., 1977: Mesoscale and convective-scale downdrafts as distinct components of squall-line structure. \textit{Mon. Wea. Rev.}, \textbf{105}, 1568--1589, doi:10.1175/1520-0493(1977)105\textless1568:MACSDA\textgreater2.0.CO;2.

\end{flushleft}

\clearpage
\section*{List of Figures}
Fig. 1. Schematic of an MCS, after Houze (2004) and Rutledge (1991).

\end{document}
